% Generated by scripts/assemble.py; corrections belong in source/corrections.json.
\examhead{40}{同济大学 2025 年数学分析试卷}{43}
% p044-b08; source PDF page 44
\par\Needspace{6\baselineskip}
\blockanchor{p044-b08}
\question{1}(10 分) 设 $\alpha, \beta$ 为正实数, 证明:\par
\nopagebreak[4]
% p044-b09; source PDF page 44
\blockanchor{p044-b09}
\displayformula{\lim _ {n \rightarrow \infty} \left(1 + \frac {\alpha + \sqrt {- 1} \beta}{n}\right) ^ {n} = e ^ {\alpha} (\cos \beta + \sqrt {- 1} \sin \beta).}
% p044-b10; source PDF page 44
\par\Needspace{6\baselineskip}
\blockanchor{p044-b10}
注：若数列 $\{a_{n}\} ,\{b_{n}\}$ 分别收敛于 $a,b,$ 则称复数列 $\{a_n + \sqrt{-1} b_n\}$ 收敛于 $a + \sqrt{-1} b,$ 记作\par
\nopagebreak[4]
% p044-b11; source PDF page 44
\blockanchor{p044-b11}
\displayformula{\lim _ {n \to \infty} (a _ {n} + \sqrt {- 1} b _ {n}) = a + \sqrt {- 1} b.}
% p044-b12; source PDF page 44
\blockanchor{p044-b12}
\question{2}(10 分) 证明: 存在 $f(x) \in C^{1}(\mathbb{R})$, 满足 $f(x)$ 在 $(- \infty, 1)$ 上恒为 0, 在 $(2, +\infty)$ 上恒为 1.\par
% p044-b13; source PDF page 44
\par\Needspace{6\baselineskip}
\blockanchor{p044-b13}
\question{3}(15 分) 设 $f$ 在 $\mathbb{R}$ 上二阶可微, 证明: 对任意的 $a, b \in \mathbb{R}, t \in (0,1)$, 存在 $\xi \in \mathbb{R}$, 满足\par
\nopagebreak[4]
% p044-b14; source PDF page 44
\blockanchor{p044-b14}
\displayformula{f ((1 - t) a + t b) = (1 - t) f (a) + t f (b) - \frac {f ^ {\prime \prime} (\xi)}{2} t (1 - t) (b - a) ^ {2}.}
% p044-b15; source PDF page 44
\blockanchor{p044-b15}
\question{4}(15 分) 设 $\alpha$ 为正实数, 讨论 $\lim_{x \to 0} \frac{\frac{1}{\sqrt{\pi}} \int_0^{+\infty} e^{-t^2} \cos(2xt) \, dt - \frac{1}{2}}{x^\alpha}$ 的敛散性及极限.\par
% p044-b16; source PDF page 44
\blockanchor{p044-b16}
\question{5}(15 分) 设 $u_{n}(x) \in R[0,1]$, 且 $\sum_{n=1}^{\infty} u_{n}(x)$ 的部分和函数列一致有界. 令 $v_{n}(x) = \int_{0}^{1} u_{n}(x) \sqrt{x + t} \, \mathrm{d}t$.\par
% p044-b17; source PDF page 44
\blockanchor{p044-b17}
(1) 证明: $v_{n}(x) \in C^{1}(0,1)$.\par
% p044-b18; source PDF page 44
\blockanchor{p044-b18}
(2) 证明: 若 $\sum_{n=1}^{\infty} v_n(x)$ 收敛于 $S(x)$, 则 $S(x) \in C^0(0,1)$.\par
% p044-b19; source PDF page 44
\blockanchor{p044-b19}
\question{6}(15 分) 设 $f, g \in C^0(\mathbb{R}^2)$, 且当 $\pmb{x} \neq \pmb{0}$ 时, 成立 $f(\pmb{x}) = g(\pmb{x}) \sin \frac{1}{\|\pmb{x}\|}$. 若 $f$ 在原点可微, 求 $\mathrm{d}f(\mathbf{0})$, 并分析 $g$ 在原点的可微性.\par
% p044-b20; source PDF page 44
\blockanchor{p044-b20}
\question{7}(15 分) 判断 $\iint_{x + y \geq 1} \frac{\sin x \sin y}{(x + y)^3} \, \mathrm{d}x \, \mathrm{d}y$ 的敛散性.\par
% p044-b21; source PDF page 44
\par\Needspace{6\baselineskip}
\blockanchor{p044-b21}
\question{8}(10 分) 曲线 $L$ 是 $x^{2} + y^{2} + z^{2} = 8$ 与 $x^{2} + y^{2} = 2z$ 的交线, 且从 $z$ 轴正向往坐标原点看去取逆时针方向. 计算\par
\nopagebreak[4]
% p044-b22; source PDF page 44
\blockanchor{p044-b22}
\displayformula{I = \int_ {L} | \sqrt {3} y - x | \mathrm{d} x + 3 y \mathrm{d} y - 2 z \mathrm{d} z.}
% p044-b23; source PDF page 44
\blockanchor{p044-b23}
\question{9}(25 分) 设 $f(x)$ 是在 $x_0$ 局部有定义的实函数, 如果 $f(x)$ 在 $x_0$ 的某邻域上可以展开为收敛于 $f(x)$ 的关于 $(x - x_0)$ 的幂级数, 则称 $f(x)$ 在 $x_0$ 点解析. 若 $f(x)$ 在开区间 $I$ 上点点解析, 则称 $f(x)$ 在 $I$ 上解析.\par
% p044-b24; source PDF page 44
\blockanchor{p044-b24}
(1) 若 $f, g$ 都为区间 $I$ 上的实解析函数, 且 $E = \{x \in I \mid f(x) = g(x)\}$ 在 $I$ 上有聚点, 证明: $f \equiv g$.\par
% p045-b00; source PDF page 45
\blockanchor{p045-b00}
(2) 设 $f, g$ 为 $\mathbb{R}$ 上的实解析函数, $w(f, g) = \det \left( \begin{array}{cc} f & g \\ f' & g' \end{array} \right)$ 为 Wronski 行列式, 证明: $w(f, g) = 0$ 当且仅当 $f, g$ 线性相关.\par
% p045-b01; source PDF page 45
\blockanchor{p045-b01}
(3) 若 $f(x), g(x) \in C^{\infty}(\mathbb{R})$, (2) 中结论是否成立？说明理由.\par
% p045-b02; source PDF page 45
\blockanchor{p045-b02}
\question{10}(20 分) 设 $P(x, y), Q(x, y)$ 为单连通区域 $\Omega$ 上的 $C^1$ 函数, 且 $\frac{\partial P}{\partial y} = \frac{\partial Q}{\partial x}$ 在 $\Omega$ 上成立. 证明: 曲线积分 $\int_{L} P \mathrm{~d}x + Q \mathrm{~d}y$ 与路径无关, 其中 $L$ 为 $\Omega$ 中的分段光滑曲线.\par
